\documentclass[11pt]{article}

\usepackage[T1]{fontenc}
\usepackage[margin=1in]{geometry}
\usepackage{amsmath,amssymb,amsthm}
\usepackage{braket}
\usepackage{array}
\usepackage{booktabs}
\usepackage{enumitem}
\usepackage[colorlinks=true,linkcolor=blue,citecolor=blue,urlcolor=blue]{hyperref}

\theoremstyle{plain}
\newtheorem{theorem}{Theorem}[section]
\newtheorem{lemma}[theorem]{Lemma}
\newtheorem{proposition}[theorem]{Proposition}
\newtheorem{corollary}[theorem]{Corollary}
\theoremstyle{definition}
\newtheorem{definition}[theorem]{Definition}
\newtheorem{remark}[theorem]{Remark}
\newtheorem{problem}[theorem]{Problem}
\newtheorem{conjecture}[theorem]{Conjecture}

\newcommand{\F}{\mathbb{F}}
\newcommand{\Z}{\mathbb{Z}}
\newcommand{\C}{\mathbb{C}}
\newcommand{\Q}{\mathbb{Q}}
\newcommand{\cat}{\mathrm{cat}}
\newcommand{\Stab}{\mathrm{Stab}}
\newcommand{\Cliff}{\mathrm{Cliff}}
\newcommand{\qq}{\mathsf{q}}
\newcommand{\bb}{\mathsf{b}}
\newcommand{\rank}{\operatorname{rank}}
\newcommand{\spn}{\operatorname{span}}
\newcommand{\Ev}{\mathcal{E}}
\newcommand{\one}{\mathbf{1}}
\newcommand{\wt}[1]{\lvert #1\rvert}
\DeclareMathOperator{\pol}{pol}

\title{\bfseries The stabilizer rank of the tensor square\\ of the six-qubit cat state}
\author{}
\date{}

\begin{document}

\maketitle

\begin{abstract}
Let $\ket{\cat_6}$ denote the six-qubit cat state built from the magic state
$\ket{T}\propto\ket{0}+e^{i\pi/4}\ket{1}$. Qassim's thesis records the bound
$\chi(\cat_6)\le 3$ and conjectures that stabilizer rank fails to be multiplicative
already for two copies, i.e.\ that $\chi(\cat_6^{\otimes2})\le 8$. This note gives a
self-contained treatment of that question. We prove
$3\le\chi(\cat_6^{\otimes2})\le 9$ with complete proofs of both bounds; we exhibit an
exact Clifford compression that identifies $\cat_6^{\otimes 2}$ with a full-support
ten-qubit state, reducing the conjecture to the single inequality
$\chi(d^{\otimes2})\le 3$ for an explicit five-qubit state $d$ of stabilizer rank two;
and we prove a family of rigidity statements that delimit where a hypothetical
eight-term decomposition can live. In particular we show that any eight-term
decomposition all of whose summands are product across the natural $6\,|\,6$ cut must
have linearly dependent left and right factor families, that the four-dimensional
``tensor-square plane'' of $d^{\otimes2}$ contains exactly four stabilizer rays and
therefore cannot support a three-term decomposition, and that the natural
maximal-isotropic enlargement of the one entangled stabilizer slice of $d^{\otimes2}$
fails a third-difference test. On the lower-bound side we prove the exact relations
$\chi(\cat_{10})\le\chi(\cat_6^{\otimes2})\le\chi(\cat_{12})$. The conjecture itself
remains open; Section~\ref{sec:open} states precisely what is missing.
All explicit identities in this note have been verified exactly in
$\mathbb{Q}(e^{i\pi/4})$; see Appendix~\ref{app:verify}.
\end{abstract}

\tableofcontents

\section{Introduction}

\subsection{Background}

Classical simulation algorithms for Clifford${}+{}$T circuits run in time controlled by
the \emph{stabilizer rank} $\chi(\psi)$: the least number of stabilizer states whose
complex span contains $\psi$. The state that matters is the magic state
$\ket{T}\propto\ket{0}+e^{i\pi/4}\ket{1}$, and the central open problem is the growth
rate of $\chi(T^{\otimes m})$. Writing
\[
\alpha:=\limsup_{m\to\infty}\frac{\log_2\chi(T^{\otimes m})}{m},
\]
the best published upper bound is $\alpha\le\log_2(3)/4\approx0.3963$, obtained by
Qassim, Pashayan and Gosset \cite{QPG2021} (see also \cite[Thm.~11]{QassimThesis}) from
the single inequality $\chi(\cat_6)\le3$ together with a contraction scheme that glues
six-qubit cat states along a tree. The exponent is exactly
$\log_2\chi(\cat_6)/(6-2)$: each vertex of the tree costs a factor $\chi(\cat_6)$ and
yields four fresh qubits of cat.

Qassim observed that this exponent is unlikely to be tight, for a structural reason:
plugging \emph{two} copies of $\cat_6$ into the same scheme costs
$\chi(\cat_6^{\otimes2})$ rather than $\chi(\cat_6)^2=9$, and equality
$\chi(\cat_6^{\otimes2})=9$ would make $\cat_6$ the first known example of a state with
multiplicative stabilizer rank. This is \cite[Eq.~(4.30)]{QassimThesis}:

\begin{conjecture}[Qassim]\label{conj:main}
$\chi(\cat_6\otimes\cat_6)\le 8$.
\end{conjecture}

\noindent Conjecture~\ref{conj:main} would give
$\alpha\le\log_2(8)/8=3/8=0.375$ \cite[Eq.~(4.31)]{QassimThesis}. The hypothesis is not
idle: Lovitz and Steffan \cite{LovitzSteffan2022} later produced the first states with
multiplicative rank, namely two-qubit states $\psi_\zeta$, depending on a parameter $\zeta\in\C^\times$, with
$\chi(\psi_\zeta)=2$ and $\chi(\psi_\zeta^{\otimes2})=4$; but their proof works only for
$\zeta$ outside a finite exceptional set, in particular for $\zeta$ transcendental over
$\Q$, and the
cat state's amplitudes are algebraic. So the question is genuinely open in both
directions.

\subsection{Results}

Throughout, $\ket{C}:=\ket{\cat_6}$ (unnormalized; see Section~\ref{sec:prelim}).
Our proved numerical result is the following.

\begin{theorem}\label{thm:main}
$3\le\chi(C\otimes C)\le9$.
\end{theorem}

\noindent Neither bound is new in substance --- the upper bound is the trivial
submultiplicative one and the lower bound follows from $\chi(\cat_6)=3$ --- but both are
proved here in full, and the proofs are the scaffolding for everything else. The
substantive content of this note is a set of exact reductions and rigidity theorems.
We single out four.

\begin{theorem}[Compression; Theorem~\ref{thm:compress}]\label{thm:introcompress}
There is a CNOT circuit $U_\Phi$ on six qubits with
$U_\Phi\ket{C}=\sqrt2\,\ket{0}\otimes\ket{\widehat C}$, where $\ket{\widehat C}$ is the
\emph{full-support} five-qubit state
$\ket{\widehat C}=\sum_{y\in\F_2^5}i^{\lceil \wt{y}/2\rceil}\ket{y}$. Consequently
$\chi(C\otimes C)=\chi(\widehat C\otimes\widehat C)$, and
$\ket{\widehat C}=2\ket{w_2}-\ket{d}$ where $\ket{w_2}=\ket{0^5}-i\ket{1^5}$ is a
stabilizer state and $\ket{d}=\sum_y i^{-\qq(y)}\ket y$ has $\chi(d)=2$. In particular
\[
\chi(C\otimes C)\;\le\;5+\chi(d\otimes d),
\]
so $\chi(d^{\otimes2})\le3$ would prove Conjecture~\ref{conj:main}.
\end{theorem}

\begin{theorem}[Product-cut obstruction; Propositions~\ref{prop:indep} and \ref{prop:sylvester}]
Suppose $C\otimes C=\sum_{j=1}^k c_j\ket{\xi_j}\otimes\ket{\eta_j}$ with all
$\xi_j,\eta_j$ six-qubit stabilizer states. If the pairwise non-proportional
$\xi_j$ are linearly independent, then $k\ge9$; the same holds with the two tensor
factors exchanged. Without any independence hypothesis, the ranks
$r=\rank(\xi_1,\dots,\xi_k)$ and $r'=\rank(\eta_1,\dots,\eta_k)$ satisfy
$r,r'\ge3$ and $r+r'\le k+1$.
\end{theorem}

\begin{theorem}[Plane rigidity; Theorems~\ref{thm:plane} and \ref{thm:squareplane}]
Let $\ket{w_0}=\sum_y\ket y$ and $\ket{w_1}=\sum_y(-1)^{\qq(y)}\ket y$. The only
stabilizer rays in $\spn\{w_0,w_1\}$ are $[w_0]$ and $[w_1]$; and the only stabilizer
rays in the four-dimensional space $\spn\{w_a\otimes w_b\}_{a,b\in\{0,1\}}$ are the four
product rays $[w_a\otimes w_b]$. Since $d\otimes d$ has all four coordinates nonzero in
that basis, no three-term decomposition of $d\otimes d$ can use only stabilizer states
from that plane.
\end{theorem}

\begin{theorem}[Two-sided sandwich; Theorem~\ref{thm:sandwich}]
$\chi(\cat_{10})\le\chi(C\otimes C)\le\chi(\cat_{12})$.
\end{theorem}

We also prove that a single entangled stabilizer slice of $d\otimes d$ does exist (the
graph $z=y+\one_5$, on which the phase of $d\otimes d$ is constant), that the natural
$64$-point maximal-isotropic enlargement of that slice is \emph{not} a stabilizer state,
and that excluding $\chi(d^{\otimes2})\le3$ reduces to two finite tasks
(Proposition~\ref{prop:reduction}).

\subsection{What is not proved}

We do not prove Conjecture~\ref{conj:main}, and we do not refute it. We also do not
improve the lower bound past $3$. Section~\ref{sec:open} states the three remaining gaps
precisely and records the search strategies that were tried and failed, so that the
reader does not repeat them.

\subsection{Organisation}

Section~\ref{sec:prelim} fixes notation and proves the elementary facts about stabilizer
rank and about phase exponents that the rest of the note uses; a reader familiar with
the affine-quadratic normal form may skim it, but the difference calculus of
Section~\ref{sec:diffcalc} is used repeatedly and is worth reading. Section~\ref{sec:three}
proves the three-term decomposition of $C$. Section~\ref{sec:compress} proves the
compression theorem. Section~\ref{sec:upper} gives the upper bound and the reduction to
$d^{\otimes2}$. Section~\ref{sec:cut} proves the product-cut obstructions.
Section~\ref{sec:rigid} contains the rigidity theorems. Section~\ref{sec:lower} proves
the lower bound and the sandwich. Section~\ref{sec:open} discusses what remains.

\section{Preliminaries}\label{sec:prelim}

\subsection{Notation}

We work over $\F_2=\{0,1\}$ and $\Z_4=\Z/4\Z$. Addition in $\F_2^n$ is written $+$ and
is bitwise; when a bit string is used inside an integer expression, its entries are
read as the integers $0$ and $1$. For $x\in\F_2^n$, $\wt{x}=\sum_j x_j\in\Z$ is the
Hamming weight. We write $\one_n=(1,1,\dots,1)\in\F_2^n$, $e_j$ for the $j$-th standard
basis vector, and $0^n$ for the zero string. An \emph{affine subspace}
$\mathcal A\subseteq\F_2^n$ is a coset $x_0+\mathcal A_0$ of a linear subspace
$\mathcal A_0$, called the \emph{direction space} of $\mathcal A$; $\dim\mathcal A:=\dim\mathcal A_0$.
An \emph{affine hyperplane} is an affine subspace of dimension $n-1$; equivalently a set
$\{x:v\cdot x=\epsilon\}$ for some $v\in\F_2^n\setminus\{0\}$ and $\epsilon\in\F_2$,
where $v\cdot x=\sum_j v_jx_j \bmod 2$.

Set $\omega:=e^{i\pi/4}$, so $\omega^2=i$ and $\omega^4=-1$. Let
\[
\Ev_n:=\{x\in\F_2^n:\wt{x}\equiv0\ (\mathrm{mod}\ 2)\}
\]
be the even-weight code, a linear subspace of dimension $n-1$. Define the two Boolean
quadratics that occur throughout:
\begin{align}
\sigma_n(x)&:=\sum_{1\le j<k\le n}x_jx_k \bmod 2 \qquad (x\in\F_2^n),\label{eq:sigma}\\
\qq(y)&:=\sigma_5(y)+\wt{y} \bmod 2 \qquad (y\in\F_2^5).\label{eq:qdef}
\end{align}
A short computation records the only properties of $\sigma_n$ we need.

\begin{lemma}\label{lem:sigma}
$\sigma_n(x)=\binom{\wt x}{2}\bmod 2$. Consequently $\sigma_6$ takes the values
$0,1,0,1$ on strings of weight $0,2,4,6$, and $\qq(y)$ depends only on $\wt y$,
taking the values $0,1,1,0,0,1$ for $\wt y=0,1,2,3,4,5$.
\end{lemma}

\begin{proof}
$\sigma_n(x)$ counts the pairs $\{j,k\}$ with $x_j=x_k=1$, of which there are
$\binom{\wt x}{2}$. The two tables follow by evaluating
$\binom{w}{2}$ and $\binom w2+w$ modulo $2$.
\end{proof}

Vectors are unnormalized unless stated otherwise; $\ket\psi$ and $\lambda\ket\psi$
($\lambda\ne0$) denote the same \emph{ray} $[\psi]$. For a vector $\ket\psi$ on $n$
qubits, its \emph{support} is $\{x\in\F_2^n:\braket{x|\psi}\ne0\}$.

\subsection{Stabilizer states, stabilizer rank, phase exponents}

\begin{definition}[Stabilizer state]\label{def:stab}
A nonzero vector $\ket\varphi\in(\C^2)^{\otimes n}$ is a \emph{stabilizer state} if
\begin{equation}\label{eq:normalform}
\ket\varphi=\kappa\sum_{x\in\mathcal A}i^{\ell(x)}(-1)^{\theta(x)}\ket x
\end{equation}
for some $\kappa\in\C^\times$, some affine subspace $\mathcal A\subseteq\F_2^n$, and,
in some affine parametrization $x=x_0+Mt$ of $\mathcal A$ by $t\in\F_2^{\delta}$
($\delta=\dim\mathcal A$, $M$ injective and $\F_2$-linear), functions
$\ell(x_0+Mt)=\sum_{j=1}^{\delta}\ell_jt_j$ with $\ell_j\in\Z_4$, and $\theta$ a quadratic
polynomial over $\F_2$ in $t$. We write $\Stab_n$ for the set of stabilizer states.
This is the standard affine-quadratic normal form; it agrees with the definition of a
stabilizer state as an element of the Clifford orbit of $\ket{0^n}$
\cite{DehaeneDeMoor2003,VandenNest2010}.
\end{definition}

It is convenient to fold the two phase factors together.

\begin{definition}[Phase exponent]\label{def:exponent}
Let $\ket\varphi$ be a nonzero vector supported on an affine subspace $\mathcal A$
whose nonzero amplitudes all have the same modulus and lie in
$\kappa\{1,i,-1,-i\}$ for some $\kappa\in\C^\times$. A \emph{phase exponent} of
$\ket\varphi$ is a function $h:\mathcal A\to\Z_4$ with
$\braket{x|\varphi}=\kappa\, i^{h(x)}$ for all $x\in\mathcal A$. It is unique once
$\kappa$ is fixed, and unique up to an additive constant otherwise; in particular all
differences of $h$ are determined by $\ket\varphi$ alone.
\end{definition}

With $h=\ell+2\theta$, \eqref{eq:normalform} reads
$\ket\varphi=\kappa\sum_{x\in\mathcal A}i^{h(x)}\ket x$. The next lemma says that the
shape ``$\Z_4$-linear plus twice quadratic'' is a property of $h$, not of the chosen
affine coordinates. We use it silently thereafter.

\begin{lemma}[Coordinate independence]\label{lem:coordfree}
Let $\mathcal A$ be an affine subspace with two affine parametrizations
$t\mapsto x_0+Mt$ and $s\mapsto x_0'+M's$, related by $t=Ns+t_0$ with $N$ invertible
over $\F_2$. If $H(t)=\sum_j\ell_jt_j+2\theta(t)$ with $\ell_j\in\Z_4$ and $\theta$
quadratic over $\F_2$, then $H(Ns+t_0)$ has the same shape in $s$.
\end{lemma}

\begin{proof}
For bits $u_1,\dots,u_N\in\{0,1\}$ one has, as integers,
$u_1\oplus\dots\oplus u_N=\sum_i u_i-2\varepsilon_2(u)+4(\cdots)$, where
$\varepsilon_2(u)=\sum_{i<i'}u_iu_{i'}$; hence modulo $4$,
$u_1\oplus\cdots\oplus u_N\equiv\sum_iu_i+2\varepsilon_2(u)$. Applying this to
$t_j=(Ns+t_0)_j$, which is an $\F_2$-sum of some of the $s_k$ and the constant
$(t_0)_j$, gives
$\ell_jt_j\equiv \ell_j\big(\text{$\Z$-linear in }s\big)+2\big(\ell_j\bmod 2\big)e_2(\cdot)\pmod 4$,
and $\varepsilon_2(\cdot)$ is a quadratic polynomial over $\F_2$. Finally an affine substitution
carries a quadratic polynomial over $\F_2$ to a quadratic polynomial, so
$2\theta(Ns+t_0)=2\theta'(s)$ with $\theta'$ quadratic.
\end{proof}

\begin{definition}[Stabilizer rank]\label{def:chi}
For a nonzero $\ket\psi\in(\C^2)^{\otimes n}$,
\[
\chi(\psi):=\min\Big\{k\ \Big|\ \ket\psi=\sum_{j=1}^k c_j\ket{\varphi_j},\
c_j\in\C,\ \ket{\varphi_j}\in\Stab_n\Big\}.
\]
\end{definition}

Since $\chi$ is defined by a span condition it depends only on the ray $[\psi]$, and
we may and do take all stabilizer states in a decomposition to be pairwise
non-proportional (proportional ones may be merged).

\begin{lemma}[Elementary operations]\label{lem:ops}
Let $\ket\psi,\ket\phi$ be nonzero states, $\lambda\in\C^\times$, $U$ a Clifford
unitary, and $\ket\tau$ a stabilizer state. Then
\begin{enumerate}[label=(\alph*)]
\item $\chi(\lambda\psi)=\chi(\psi)$ and $\chi(U\psi)=\chi(\psi)$;
\item $\chi(\psi\otimes\tau)=\chi(\psi)$;
\item $\chi(\psi\otimes\phi)\le\chi(\psi)\chi(\phi)$;
\item if $P$ is a Pauli operator with $P^2=I$ and $\Pi=\tfrac12(I+P)$, then
$\Pi\ket\psi=0$ or $\chi(\Pi\psi)\le\chi(\psi)$;
\item if $\ket\tau$ is a stabilizer state on a subset $S$ of the qubits and
$\ket{\psi'}:=(\mathrm{id}_{S^c}\otimes\bra\tau)\ket\psi$ is nonzero, then
$\chi(\psi')\le\chi(\psi)$.
\end{enumerate}
\end{lemma}

\begin{proof}
(a) Cliffords permute $\Stab_n$ and scaling does not change spans.

(c) If $\psi=\sum_ia_i\varphi_i$ and $\phi=\sum_jb_j\varphi_j'$ then
$\psi\otimes\phi=\sum_{i,j}a_ib_j\,\varphi_i\otimes\varphi_j'$, and a tensor product of
stabilizer states is a stabilizer state: its support is the Cartesian product
$\mathcal A\times\mathcal A'$, again affine, and in the concatenated affine coordinates
its phase exponent is $h(t)+h'(t')$, again of the shape of
Definition~\ref{def:stab}.

(b) ``$\le$'' is (c) with $\chi(\tau)=1$. For ``$\ge$'', let
$\psi\otimes\tau=\sum_jc_j\varphi_j$ be a minimal decomposition and contract the second
factor with $\bra\tau$: since $\braket{\tau|\tau}\ne0$ we get
$\braket{\tau|\tau}\ket\psi=\sum_jc_j(\mathrm{id}\otimes\bra\tau)\ket{\varphi_j}$, and by
(e) below each term is zero or a multiple of a stabilizer state.

(d) Every Pauli $P$ with $P^2=I$ satisfies $P=\pm U Z_1U^\dagger$ for some Clifford $U$,
so by (a) we may assume $\Pi=\frac12(I\pm Z_1)$, i.e.\ $\Pi$ restricts the support to
the affine hyperplane $\{x_1=\epsilon\}$. If $\ket\varphi$ is a stabilizer state with
support $\mathcal A$ and phase exponent $h$, then $\Pi\ket\varphi$ has support
$\mathcal A\cap\{x_1=\epsilon\}$, which is empty or affine, and phase exponent the
restriction of $h$, which retains the shape of Definition~\ref{def:stab} by
Lemma~\ref{lem:coordfree} applied to any affine parametrization of the intersection.
Applying $\Pi$ to a $\chi(\psi)$-term decomposition of $\psi$ therefore yields at most
$\chi(\psi)$ terms.

(e) Write $\ket\tau=W\ket{0^{|S|}}$ for a Clifford $W$ on $S$; then
$\bra\tau=\bra{0^{|S|}}W^\dagger$, so by (a) we may assume $\ket\tau=\ket{0^{|S|}}$.
Contraction with $\bra{0^{|S|}}$ is the composition of the Pauli projections
$\frac12(I+Z_j)$, $j\in S$, of (d) with the deletion of the (now product) qubits in
$S$, and deletion of a tensor factor equal to $\ket{0}$ does not change stabilizer
rank by (b).
\end{proof}

\subsection{A difference calculus for phase exponents}\label{sec:diffcalc}

The single tool used to \emph{exclude} stabilizer states below is the following
elementary calculus. For $a\in\mathcal A_0$ define the difference operator
\[
(\Delta_a h)(x):=h(x+a)-h(x)\in\Z_4,\qquad x\in\mathcal A,
\]
and iterate, so that
$(\Delta_a\Delta_bh)(x)=h(x+a+b)-h(x+a)-h(x+b)+h(x)$.

For a quadratic polynomial $\theta$ over $\F_2$ let
$\pol_\theta(a,b):=\theta(a+b)+\theta(a)+\theta(b)+\theta(0)$ denote its
\emph{polar form}; it is symmetric, $\F_2$-bilinear, alternating
($\pol_\theta(a,a)=0$), and independent of any base point.

\begin{lemma}[Difference calculus]\label{lem:diff}
Let $\ket\varphi$ be a stabilizer state with support $\mathcal A$, direction space
$\mathcal A_0$ and phase exponent $h$. Then there is a symmetric $\F_2$-bilinear form
$\Lambda$ on $\mathcal A_0$ such that for all $x\in\mathcal A$ and all
$a,b\in\mathcal A_0$,
\begin{equation}\label{eq:secondiff}
(\Delta_a\Delta_bh)(x)\equiv 2\,\Lambda(a,b) \pmod 4 .
\end{equation}
In particular:
\begin{enumerate}[label=(\roman*)]
\item \textbf{(parity)} every mixed second difference of $h$ is even;
\item \textbf{(base-point freedom)} every mixed second difference of $h$ is independent
of the base point $x$; equivalently every third difference
$\Delta_a\Delta_b\Delta_ch$ vanishes in $\Z_4$;
\item \textbf{(reduction mod $2$)} $h\bmod2$ is an $\F_2$-affine function on $\mathcal A$.
\end{enumerate}
\end{lemma}

\begin{proof}
Choose affine coordinates $x=x_0+Mt$ as in Definition~\ref{def:stab}, so that
$H(t):=h(x_0+Mt)=\sum_j\ell_jt_j+2\theta(t)$; by Lemma~\ref{lem:coordfree} nothing
depends on the choice. Let $a=Ma'$, $b=Mb'$.

For a single coordinate, the map $t_j\mapsto \ell_jt_j$ has second difference
\[
\ell_j\big[(t_j\oplus a_j'\oplus b_j')-(t_j\oplus a_j')-(t_j\oplus b_j')+t_j\big].
\]
The bracket vanishes unless $a_j'=b_j'=1$, in which case it equals
$t_j-2(1-t_j)+t_j=4t_j-2\equiv2\pmod4$. Hence
\[
\Delta_{a'}\Delta_{b'}\Big(\sum_j\ell_jt_j\Big)\equiv 2\sum_j\ell_ja_j'b_j'\pmod4 ,
\]
which is independent of $t$. For the quadratic part,
\[
\Delta_{a'}\Delta_{b'}\big(2\theta\big)(t)=2\big[\theta(t+a'+b')-\theta(t+a')-\theta(t+b')+\theta(t)\big]
\equiv 2\pol_\theta(a',b')\pmod 4 ,
\]
again independent of $t$, since the bracket reduces modulo $2$ to the polar form of a
quadratic. Adding the two contributions gives \eqref{eq:secondiff} with
$\Lambda(a,b)=\sum_j\bar\ell_ja_j'b_j'+\pol_\theta(a',b')$ where
$\bar\ell_j=\ell_j\bmod2$; this is symmetric and $\F_2$-bilinear.

(i) and (ii) are immediate from \eqref{eq:secondiff}. For (iii), reduce
$H=\sum_j\ell_jt_j+2\theta$ modulo $2$ to get $\sum_j\bar\ell_jt_j$, an $\F_2$-linear
function of $t$, hence an $\F_2$-affine function of $x$ on $\mathcal A$.
\end{proof}

\begin{corollary}[Non-stabilizer test]\label{cor:test}
Let $\ket\psi$ be supported on an affine subspace $\mathcal A$ with all amplitudes of
equal modulus, and let $h$ be a phase exponent. If either
\begin{enumerate}[label=(\alph*)]
\item some mixed second difference $(\Delta_a\Delta_bh)(x)$ is odd, or
\item some mixed second difference $(\Delta_a\Delta_bh)(x)$ depends on $x$,
\end{enumerate}
then $\ket\psi$ is not a stabilizer state.
\end{corollary}

\begin{proof}
Immediate from Lemma~\ref{lem:diff}(i)--(ii); the differences of $h$ do not depend on
the choice of $h$ (Definition~\ref{def:exponent}).
\end{proof}

Finally we record the isotropy bound used in Section~\ref{sec:rigid}.

\begin{lemma}[Isotropy bound]\label{lem:isotropy}
Let $\bb$ be an alternating $\F_2$-bilinear form on $V=\F_2^N$ with radical
$\mathcal R=\{v:\bb(v,\cdot)=0\}$ of dimension $\rho$. Then every totally isotropic
subspace $\mathcal W\subseteq V$ (i.e.\ $\bb|_{\mathcal W\times\mathcal W}=0$) satisfies
\[
\dim\mathcal W\le \rho+\tfrac{N-\rho}{2}.
\]
\end{lemma}

\begin{proof}
$\bb$ descends to a nondegenerate alternating form on $V/\mathcal R$, of even dimension
$N-\rho$, in which a totally isotropic subspace has dimension at most $(N-\rho)/2$.
The image of $\mathcal W$ in $V/\mathcal R$ is totally isotropic, and
$\dim\mathcal W\le\dim(\mathcal W\cap\mathcal R)+\dim(\mathcal W+\mathcal R)/\mathcal R
\le\rho+(N-\rho)/2$.
\end{proof}

\subsection{Magic states and cat states}

\begin{definition}\label{def:T}
Put $\ket T:=\ket0+\omega\ket1$ and $\ket{T^\perp}:=\ket0-\omega\ket1=Z\ket T$; these
are unnormalized and orthogonal. For $n\ge1$,
\[
\ket{\cat_n}:=\tfrac1{\sqrt2}\big(\ket T^{\otimes n}+\ket{T^\perp}^{\otimes n}\big),
\qquad \ket C:=\ket{\cat_6}.
\]
\end{definition}

Normalization conventions differ across the literature: \cite{QPG2021} uses the unit
vector $2^{-1/2}(\ket0+\omega\ket1)$, so that its cat state is a fixed nonzero multiple
of ours. By Lemma~\ref{lem:ops}(a) this is immaterial for stabilizer rank, and we use
the cited values without further comment.

\begin{lemma}[Amplitudes of the cat state]\label{lem:catamp}
\begin{equation}\label{eq:catamp}
\ket C=\sqrt2\sum_{x\in\Ev_6}\omega^{\wt x}\ket x .
\end{equation}
More generally $\ket{\cat_n}=\sqrt2\sum_{x\in\Ev_n}\omega^{\wt x}\ket x$.
\end{lemma}

\begin{proof}
$\ket T^{\otimes n}=\sum_{x\in\F_2^n}\omega^{\wt x}\ket x$ and
$\ket{T^\perp}^{\otimes n}=\sum_x(-1)^{\wt x}\omega^{\wt x}\ket x$, so the sum has
amplitude $\omega^{\wt x}(1+(-1)^{\wt x})$, which is $2\omega^{\wt x}$ for even
$\wt x$ and $0$ otherwise. Divide by $\sqrt2$.
\end{proof}

We shall use the following two facts from the literature.

\begin{theorem}[{\cite[\S4.2]{QassimThesis}, \cite{QPG2021}}]\label{thm:cited}
$\chi(\cat_6)=3$. The upper bound $\chi(\cat_6)\le3$ is
\cite[\S4.2]{QassimThesis} and is reproved from scratch in
Theorem~\ref{thm:threeterm} below. The lower bound is proved in the appendix of
\cite{QPG2021} in the form $\chi(\cat_6)\ge\chi(\cat_5)\ge3$, by classifying pairs of
stabilizer states up to a Clifford into a canonical form and comparing Pauli spectra;
we use it as a citation and do not reproduce that finite case analysis.
\end{theorem}

\section{An exact three-term decomposition of the cat state}\label{sec:three}

Define three vectors on six qubits:
\begin{equation}\label{eq:phis}
\ket{\varphi_1}:=\ket{0^6}-i\ket{1^6},\qquad
\ket{\varphi_2}:=\sum_{x\in\Ev_6}\ket x,\qquad
\ket{\varphi_3}:=\sum_{x\in\Ev_6}(-1)^{\sigma_6(x)}\ket x ,
\end{equation}
with $\sigma_6$ as in \eqref{eq:sigma}, and three scalars
\begin{equation}\label{eq:as}
a_1:=2\sqrt2,\qquad a_2:=\frac{-1+i}{\sqrt2},\qquad a_3:=\frac{-1-i}{\sqrt2}.
\end{equation}

\begin{lemma}\label{lem:phistab}
$\ket{\varphi_1},\ket{\varphi_2},\ket{\varphi_3}$ are stabilizer states.
\end{lemma}

\begin{proof}
For $\ket{\varphi_1}$: the support is the repetition code
$\mathcal A=\{0^6,\one_6\}$, a linear subspace of dimension $1$, parametrized by
$x=t\one_6$, $t\in\F_2$. The amplitudes are $1$ and $-i=i^3$, so
$h(t)=3t$ is $\Z_4$-linear and \eqref{eq:normalform} holds with $\theta=0$,
$\kappa=1$.

For $\ket{\varphi_2}$ and $\ket{\varphi_3}$: the support is $\Ev_6$, a linear subspace
of dimension $5$, parametrized by $t\in\F_2^5\mapsto x=(t_1+\dots+t_5,t_1,\dots,t_5)$.
The phase exponents are $h=0$ and $h=2\sigma_6$ respectively, and $\sigma_6$ restricted
to $\Ev_6$ is a quadratic polynomial in $t$ (substituting an affine expression into a
quadratic gives a quadratic).
\end{proof}

\begin{theorem}[Three-term decomposition]\label{thm:threeterm}
\begin{equation}\label{eq:threeterm}
\ket C=a_1\ket{\varphi_1}+a_2\ket{\varphi_2}+a_3\ket{\varphi_3}
=2\sqrt2\,\ket{\varphi_1}+\frac{-1+i}{\sqrt2}\ket{\varphi_2}+\frac{-1-i}{\sqrt2}\ket{\varphi_3}.
\end{equation}
In particular $\chi(C)\le3$, and by Theorem~\ref{thm:cited}, $\chi(C)=3$.
\end{theorem}

\begin{proof}
Both sides vanish on odd-weight strings by \eqref{eq:catamp} and \eqref{eq:phis}, so it
suffices to compare amplitudes at $x\in\Ev_6$. By Lemma~\ref{lem:sigma},
$(-1)^{\sigma_6(x)}$ equals $+1,-1,+1,-1$ according as $\wt x=0,2,4,6$; and
$\ket{\varphi_1}$ contributes only at $\wt x\in\{0,6\}$. Using
$a_2+a_3=-\sqrt2$ and $a_2-a_3=i\sqrt2$, the right-hand side has amplitudes
\[
\begin{array}{c|c|c}
\wt x & \text{amplitude of }a_1\varphi_1+a_2\varphi_2+a_3\varphi_3 & \sqrt2\,\omega^{\wt x}\\ \hline
0 & a_1+a_2+a_3=2\sqrt2-\sqrt2=\sqrt2 & \sqrt2\\
2 & a_2-a_3=i\sqrt2 & \sqrt2\,\omega^2=i\sqrt2\\
4 & a_2+a_3=-\sqrt2 & \sqrt2\,\omega^4=-\sqrt2\\
6 & -ia_1+a_2-a_3=-2\sqrt2\,i+\sqrt2\,i=-i\sqrt2 & \sqrt2\,\omega^6=-i\sqrt2 .
\end{array}
\]
The two columns agree, which by \eqref{eq:catamp} proves \eqref{eq:threeterm}.
\end{proof}

\begin{lemma}\label{lem:phiindep}
$\ket{\varphi_1},\ket{\varphi_2},\ket{\varphi_3}$ are linearly independent.
\end{lemma}

\begin{proof}
Suppose $\lambda_1\varphi_1+\lambda_2\varphi_2+\lambda_3\varphi_3=0$. Evaluating at a
string of weight $2$ gives $\lambda_2-\lambda_3=0$ and at a string of weight $4$ gives
$\lambda_2+\lambda_3=0$, so $\lambda_2=\lambda_3=0$; then evaluating at $0^6$ gives
$\lambda_1=0$.
\end{proof}

\section{Clifford compression to five qubits}\label{sec:compress}

The cat state is supported on the even-weight code, which is a hyperplane; a CNOT
circuit therefore compresses it to a full-support state on five qubits. Making this
explicit removes a redundant qubit from every statement below and turns the target of
Conjecture~\ref{conj:main} into a full-support ten-qubit problem.

Let $\Phi:\F_2^6\to\F_2^6$ be the invertible $\F_2$-linear map
\begin{equation}\label{eq:Phi}
\Phi(x_1,\dots,x_6)=(p,y_1,\dots,y_5):=(x_1+\dots+x_6,\ x_2,\dots,x_6),
\end{equation}
with inverse $\Phi^{-1}(p,y)=(p+y_1+\dots+y_5,y_1,\dots,y_5)$, and let $U_\Phi$ be the
unitary permuting the computational basis by $\ket x\mapsto\ket{\Phi(x)}$. Since $\Phi$
is $\F_2$-linear and invertible, $U_\Phi$ is a product of CNOT gates, hence Clifford.
Note $x\in\Ev_6\iff p=0$.

\begin{lemma}\label{lem:sigmaq}
For $x\in\Ev_6$ with $\Phi(x)=(0,y)$ we have $\sigma_6(x)=\qq(y)$, with $\qq$ as in
\eqref{eq:qdef}.
\end{lemma}

\begin{proof}
Write $n:=\wt y\in\{0,\dots,5\}$. On $\Ev_6$ we have $x_1=y_1+\dots+y_5=n\bmod2$ and
$x_{j+1}=y_j$, so $\wt x=n+(n\bmod 2)$. By Lemma~\ref{lem:sigma},
$\sigma_6(x)=\binom{\wt x}2\bmod2$ and $\qq(y)=\binom n2+n\bmod2$. If $n$ is even then
$\wt x=n$ and both sides equal $\binom n2$. If $n$ is odd then $\wt x=n+1$ and
$\binom{n+1}2=\binom n2+n$, as required.
\end{proof}

Define, on five qubits,
\begin{align}
\ket{w_0}&:=\sum_{y\in\F_2^5}\ket y, &
\ket{w_1}&:=\sum_{y\in\F_2^5}(-1)^{\qq(y)}\ket y, &
\ket{w_2}&:=\ket{0^5}-i\ket{1^5},\label{eq:ws}\\
\ket d&:=\sum_{y\in\F_2^5}i^{-\qq(y)}\ket y, &
\ket{\widehat C}&:=\sum_{y\in\F_2^5}i^{\lceil \wt y/2\rceil}\ket y. & &
\label{eq:dC}
\end{align}
Here $\qq(y)\in\{0,1\}$ is read as an integer in the exponent. The vectors
$\ket{w_0},\ket{w_1},\ket{w_2}$ are stabilizer states, by the same argument as
Lemma~\ref{lem:phistab}: $\ket{w_0},\ket{w_1}$ have full (hence affine) support with
phase exponents $0$ and $2\qq$, and $\ket{w_2}$ has support $\{0^5,\one_5\}$ with
exponent $3t$.

\begin{theorem}[Compression]\label{thm:compress}
\begin{equation}\label{eq:compress}
U_\Phi\ket{\varphi_1}=\ket0\otimes\ket{w_2},\quad
U_\Phi\ket{\varphi_2}=\ket0\otimes\ket{w_0},\quad
U_\Phi\ket{\varphi_3}=\ket0\otimes\ket{w_1},\quad
U_\Phi\ket{C}=\sqrt2\,\ket0\otimes\ket{\widehat C},
\end{equation}
and
\begin{equation}\label{eq:Chat}
\ket{\widehat C}=2\ket{w_2}-\ket d,\qquad
\ket d=\frac{1-i}{2}\big(\ket{w_0}+i\ket{w_1}\big).
\end{equation}
Consequently
\begin{equation}\label{eq:chieq}
\chi(C)=\chi(\widehat C)=3,\qquad
\chi(C\otimes C)=\chi(\widehat C\otimes\widehat C).
\end{equation}
\end{theorem}

\begin{proof}
The first three identities in \eqref{eq:compress} are immediate:
$\Phi(0^6)=(0,0^5)$ and $\Phi(\one_6)=(0,\one_5)$, giving $\ket{w_2}$; and $\Phi$ maps
$\Ev_6$ bijectively onto $\{0\}\times\F_2^5$, carrying the phase $1$ to $1$ and, by
Lemma~\ref{lem:sigmaq}, the phase $(-1)^{\sigma_6}$ to $(-1)^{\qq}$.

For the fourth, let $x\in\Ev_6$ with $\Phi(x)=(0,y)$ and $n=\wt y$. As in the proof of
Lemma~\ref{lem:sigmaq}, $\wt x=n+(n\bmod2)=2\lceil n/2\rceil$, so by \eqref{eq:catamp}
the amplitude of $\ket C$ at $x$ is
$\sqrt2\,\omega^{2\lceil n/2\rceil}=\sqrt2\,i^{\lceil n/2\rceil}$, which is
$\sqrt2$ times the amplitude of $\ket{\widehat C}$ at $y$.

For \eqref{eq:Chat}: the amplitude of $\ket d$ is $1$ where $\qq=0$ and $i^{-1}=-i$
where $\qq=1$; the vector $\frac{1-i}2(w_0+iw_1)$ has amplitude
$\frac{1-i}2(1+i)=1$ where $\qq=0$ and $\frac{1-i}{2}(1-i)=-i$ where $\qq=1$. For the
first identity of \eqref{eq:Chat}, compare $2w_2-d$ with $\widehat C$ at each $y$,
using Lemma~\ref{lem:sigma} for the values of $\qq$:
for $y\notin\{0^5,\one_5\}$ we need $i^{\lceil n/2\rceil}=-i^{-\qq(y)}$ with $n=\wt y$, i.e.\
$\lceil n/2\rceil\equiv2-\qq(y)\pmod4$, which holds since
$(\lceil n/2\rceil,\qq)=(1,1),(1,1),(2,0),(2,0)$ for $n=1,2,3,4$; at $y=0^5$ the left
side is $i^0=1$ and the right side is $2-1=1$; at $y=\one_5$ the left side is
$i^3=-i$ and the right side is $-2i-i^{-1}=-2i+i=-i$.

Finally \eqref{eq:chieq} follows from Lemma~\ref{lem:ops}(a),(b): $U_\Phi$ is Clifford
and $\ket 0$ is a stabilizer state, so
$\chi(C)=\chi(U_\Phi C)=\chi(\ket0\otimes\widehat C)=\chi(\widehat C)$; the value $3$
is Theorem~\ref{thm:cited} together with Theorem~\ref{thm:threeterm}. Applying
$U_\Phi\otimes U_\Phi$ and deleting the two ancillary qubits gives the second identity.
\end{proof}

\begin{remark}
Combining \eqref{eq:compress} with \eqref{eq:threeterm} recovers the three-term
decomposition in compressed form:
$\ket{\widehat C}=2\ket{w_2}+\frac{-1+i}{2}\ket{w_0}+\frac{-1-i}{2}\ket{w_1}$, which is
\eqref{eq:Chat} expanded.
\end{remark}

\section{The upper bound, and reduction to a rank-two tensor square}\label{sec:upper}

\begin{theorem}[Nine-term bound]\label{thm:nine}
\begin{equation}\label{eq:nine}
\ket C\otimes\ket C=\sum_{u,v\in\{1,2,3\}}a_ua_v\,\ket{\varphi_u}\otimes\ket{\varphi_v},
\end{equation}
an exact decomposition into nine linearly independent stabilizer states. Hence
$\chi(C\otimes C)\le9$.
\end{theorem}

\begin{proof}
Expand \eqref{eq:threeterm} in both factors. Each $\varphi_u\otimes\varphi_v$ is a
stabilizer state by the argument in Lemma~\ref{lem:ops}(c). Independence: by
Lemma~\ref{lem:phiindep} the family $\{\varphi_u\}$ is independent, and the set of
pairwise tensor products of two independent families is independent.
\end{proof}

Independence of the nine summands says only that no term of \eqref{eq:nine} can be
deleted; it is not an obstruction to an eight-term decomposition using different, and
possibly entangled, stabilizer states. The next reduction isolates where such a saving
could come from.

\begin{theorem}[Reduction to $d^{\otimes2}$]\label{thm:reduce}
With $\ket d$ as in \eqref{eq:dC},
\begin{equation}\label{eq:reduce}
\chi(C\otimes C)\ \le\ 5+\chi(d\otimes d).
\end{equation}
In particular $\chi(d^{\otimes2})\le3$ implies Conjecture~\ref{conj:main}, and
$\chi(d^{\otimes2})=4$ recovers only the bound $9$.
\end{theorem}

\begin{proof}
By \eqref{eq:chieq} it suffices to bound $\chi(\widehat C\otimes\widehat C)$. Using
$\ket{\widehat C}=2\ket{w_2}-\ket d$,
\[
\ket{\widehat C}^{\otimes2}
=4\,\ket{w_2}\otimes\ket{w_2}
-2\big(\ket{w_2}\otimes\ket d+\ket d\otimes\ket{w_2}\big)
+\ket d\otimes\ket d .
\]
The first term is a single stabilizer state. Expanding $\ket d$ by \eqref{eq:Chat}
turns each of the two mixed terms into two stabilizer states, e.g.
$\ket{w_2}\otimes\ket d=\frac{1-i}2\big(\ket{w_2}\otimes\ket{w_0}
+i\,\ket{w_2}\otimes\ket{w_1}\big)$. This accounts for $1+2+2=5$ stabilizer states,
and the last term costs $\chi(d\otimes d)$.
\end{proof}

\begin{lemma}\label{lem:chid}
$\chi(d)=2$.
\end{lemma}

\begin{proof}
$\chi(d)\le2$ by \eqref{eq:Chat}. For $\chi(d)\ge2$ we must show $\ket d$ is not a
stabilizer state. It has full support $\F_2^5$ and phase exponent $h(y)=-\qq(y)\bmod 4$.
By Lemma~\ref{lem:sigma}, $\qq(0)=0$, $\qq(e_1)=\qq(e_2)=1$ (weight $1$) and
$\qq(e_1+e_2)=1$ (weight $2$), so
\[
(\Delta_{e_1}\Delta_{e_2}h)(0)=-\qq(e_1+e_2)+\qq(e_1)+\qq(e_2)-\qq(0)=-1+1+1-0=1,
\]
which is odd. Corollary~\ref{cor:test}(a) applies.
\end{proof}

Thus Conjecture~\ref{conj:main} would follow from the multiplicativity failure
$\chi(d^{\otimes2})\le3$ for an explicit rank-two five-qubit state. This is exactly the
regime in which Lovitz and Steffan \cite[Thm.~4.1]{LovitzSteffan2022} proved
multiplicativity \emph{does} hold --- but only for two-qubit states $\psi_\zeta$ and
only for $\zeta$ outside a finite exceptional set, in particular for $\zeta$
transcendental over $\Q$. The amplitudes of $\ket d$ lie in $\Q(i)$, so their theorem
gives no information here.

\section{Obstructions for decompositions that are product across the cut}\label{sec:cut}

Write $C\otimes C$ as a state on $6+6$ qubits and call the induced bipartition the
\emph{$6\,|\,6$ cut}. This section shows that a hypothetical eight-term decomposition
cannot consist of summands that are products across that cut in any ``generic'' way.
Both statements are unconditional but neither constrains stabilizer states entangled
across the cut, which are allowed in Definition~\ref{def:chi}; this is the reason they
fall short of refuting Conjecture~\ref{conj:main}.

\begin{proposition}[Independent factors force nine terms]\label{prop:indep}
Suppose
\begin{equation}\label{eq:prodcut}
\ket C\otimes\ket C=\sum_{j=1}^kc_j\,\ket{\xi_j}\otimes\ket{\eta_j},
\qquad \ket{\xi_j},\ket{\eta_j}\in\Stab_6,\ c_j\in\C^\times .
\end{equation}
Group the terms by proportionality of the left factors and let
$\ket{\xi_{(1)}},\dots,\ket{\xi_{(m)}}$ be the resulting pairwise
non-proportional left factors. If $\ket{\xi_{(1)}},\dots,\ket{\xi_{(m)}}$ are
linearly independent, then $k\ge9$. The same conclusion holds with the roles of the two
tensor factors exchanged.
\end{proposition}

\begin{proof}
Let $G_p\subseteq\{1,\dots,k\}$ be the group of indices whose left factor is
proportional to $\ket{\xi_{(p)}}$; absorbing the proportionality constants into the
coefficients, \eqref{eq:prodcut} becomes
\begin{equation}\label{eq:grouped}
\ket C\otimes\ket C=\sum_{p=1}^m \ket{\xi_{(p)}}\otimes\ket{u_p},
\qquad \ket{u_p}=\sum_{j\in G_p}c_j'\ket{\eta_j}.
\end{equation}
Let $f$ be any linear functional on the second factor. Applying
$\mathrm{id}\otimes f$ to \eqref{eq:grouped} gives
$f(C)\ket C=\sum_p f(u_p)\ket{\xi_{(p)}}$. Choosing one $f$ with $f(C)\ne0$ shows
that $\ket C\in\spn\{\xi_{(p)}\}$, so by independence there are unique scalars
$\lambda_p$ with $\ket C=\sum_p\lambda_p\ket{\xi_{(p)}}$. Substituting back and
comparing coordinates in the independent family $\{\xi_{(p)}\}$ gives
$f(u_p)=\lambda_pf(C)$ for every $f$ and every $p$, i.e.
\[
\ket{u_p}=\lambda_p\ket C\qquad(1\le p\le m).
\]
Now $\ket C=\sum_p\lambda_p\ket{\xi_{(p)}}$ expresses $C$ in pairwise
non-proportional stabilizer states, so by $\chi(C)=3$ (Theorem~\ref{thm:cited}) at
least three of the $\lambda_p$ are nonzero. For each such $p$, the identity
$\ket{u_p}=\lambda_p\ket C$ expresses a nonzero multiple of $C$ as a combination of the
$|G_p|$ stabilizer states $\{\eta_j\}_{j\in G_p}$, whence $|G_p|\ge\chi(C)=3$.
Therefore $k=\sum_p|G_p|\ge3\cdot3=9$.
\end{proof}

\begin{proposition}[Rank inequality]\label{prop:sylvester}
In the situation \eqref{eq:prodcut}, let $\mathsf X$ (resp.\ $\mathsf Y$) be the
$64\times k$ matrix whose columns are the amplitude vectors of the $\xi_j$ (resp.\
$\eta_j$), and put $r=\rank\mathsf X$, $r'=\rank\mathsf Y$. Then
\[
r\ge3,\qquad r'\ge3,\qquad r+r'\le k+1 .
\]
In particular an eight-term decomposition into product stabilizer states would require
$3\le r,r'$ and $r+r'\le9$, so at least one of the two families is linearly dependent.
\end{proposition}

\begin{proof}
Matricize $C\otimes C$ across the cut: let $M\in\C^{64\times64}$ have entries
$M_{x,x'}=\braket{x|C}\braket{x'|C}$, so $M=\mathbf c\,\mathbf c^{\mathsf T}$ where
$\mathbf c$ is the amplitude vector of $C$; thus $\rank M=1$. Equation
\eqref{eq:prodcut} reads $M=\mathsf X\,\mathsf D\,\mathsf Y^{\mathsf T}$ with
$\mathsf D=\operatorname{diag}(c_1,\dots,c_k)$ invertible. Sylvester's rank inequality
applied to the product $(\mathsf X\mathsf D)\cdot\mathsf Y^{\mathsf T}$, whose inner
dimension is $k$, gives
\[
1=\rank M\ \ge\ \rank(\mathsf X\mathsf D)+\rank(\mathsf Y^{\mathsf T})-k=r+r'-k .
\]
For $r\ge3$: the column space of $M$ is $\C\mathbf c$ and is contained in the column
space of $\mathsf X$, so $\ket C\in\spn\{\xi_j\}$. If $r\le2$, choose $r$ linearly
independent columns of $\mathsf X$; they span $\spn\{\xi_j\}\ni\ket C$, exhibiting
$C$ as a combination of at most two stabilizer states and contradicting $\chi(C)=3$.
The argument for $r'$ is symmetric.
\end{proof}

\section{Rigidity of the residual state and its tensor square}\label{sec:rigid}

Throughout this section we work on five and ten qubits with $\qq$ as in
\eqref{eq:qdef} and $w_0,w_1,d$ as in \eqref{eq:ws}--\eqref{eq:dC}. We write
$\bb$ for the polar form of $\qq$.

\begin{lemma}[Structure of $\qq$]\label{lem:qstruct}
\begin{enumerate}[label=(\alph*)]
\item $\bb(y,z)=\big(\sum_jy_j\big)\big(\sum_jz_j\big)+y\cdot z \pmod 2$, an
alternating $\F_2$-bilinear form whose radical is $\{0^5,\one_5\}$, of dimension $1$.
\item $\qq(y+\one_5)=\qq(y)+1$; hence each level set $\qq^{-1}(0)$, $\qq^{-1}(1)$ has
exactly $16$ elements.
\item Neither level set of $\qq$ is an affine subspace of $\F_2^5$.
\end{enumerate}
\end{lemma}

\begin{proof}
(a) The linear part $\wt y$ of $\qq$ does not contribute to the polar form, so
$\bb=\pol_{\sigma_5}$. Expanding,
$\sigma_5(y+z)=\sigma_5(y)+\sigma_5(z)+\sum_{j\ne k}y_jz_k$ and
$\sum_{j\ne k}y_jz_k=(\sum_jy_j)(\sum_kz_k)-\sum_jy_jz_j$, which modulo $2$ is the
stated formula. It is alternating since $\bb(y,y)=\wt y^2+\wt y\equiv0$. For the
radical, $\bb(y,e_j)=\big(\sum_ky_k\big)+y_j$, so $y$ is in the radical iff
$y_j=\sum_ky_k$ for every $j$, i.e.\ iff all $y_j$ are equal; both $0^5$ and $\one_5$
satisfy this (for $\one_5$, $\sum_k1=5\equiv1$).

(b) By Lemma~\ref{lem:sigma}, $\qq(y)$ depends only on $n=\wt y$ through the table
$(0,1,1,0,0,1)$ for $n=0,\dots,5$; and $\wt{y+\one_5}=5-n$. The table read backwards is
$(1,0,0,1,1,0)$, the complement, which is the assertion. It follows that
$y\mapsto y+\one_5$ is a bijection between the two level sets, so each has $32/2=16$
elements.

(c) $\qq^{-1}(0)$ contains $0^5$, so if it were affine it would be a linear subspace.
It contains $u_1=e_1+e_2+e_3$ and $u_2=e_1+e_2+e_4$ (weight $3$, where $\qq=0$) but not
$u_1+u_2=e_3+e_4$ (weight $2$, where $\qq=1$); so it is not closed under addition. If
$\qq^{-1}(1)$ were affine then, having $16=2^4$ elements, it would be an affine
hyperplane, and its complement $\qq^{-1}(0)$ would be the opposite coset of the same
hyperplane, hence affine --- contradiction.
\end{proof}

\subsection{The plane spanned by $w_0$ and $w_1$}

\begin{theorem}[Plane rigidity]\label{thm:plane}
The only stabilizer rays contained in $\spn\{\ket{w_0},\ket{w_1}\}$ are $[w_0]$ and
$[w_1]$.
\end{theorem}

\begin{proof}
A vector $\ket\psi=\mu_0\ket{w_0}+\mu_1\ket{w_1}$ has amplitude $\mu_0+\mu_1$ on
$\qq^{-1}(0)$ and $\mu_0-\mu_1$ on $\qq^{-1}(1)$; in particular its amplitude is
constant on each level set.

\emph{Case 1: both amplitudes nonzero.} Then $\ket\psi$ has full support, so if it is a
stabilizer state all its amplitudes have equal modulus and their pairwise ratios are
powers of $i$. Hence $(\mu_0-\mu_1)/(\mu_0+\mu_1)\in\{1,-1,i,-i\}$. The value $1$ forces
$\mu_1=0$, giving $[w_0]$; the value $-1$ forces $\mu_0=0$, giving $[w_1]$. The values
$\pm i$ give $\ket\psi\propto\sum_yi^{\pm\qq(y)}\ket y$, whose phase exponent is
$h=\pm\qq$; exactly as in Lemma~\ref{lem:chid},
$(\Delta_{e_1}\Delta_{e_2}h)(0)=\pm1$ is odd, so Corollary~\ref{cor:test}(a) rules
these out.

\emph{Case 2: one amplitude vanishes.} Then the support of $\ket\psi$ is a level set of
$\qq$, which is not affine by Lemma~\ref{lem:qstruct}(c); but the support of a
stabilizer state is affine.
\end{proof}

\subsection{The tensor-square plane}

The four product states $\ket{w_a}\otimes\ket{w_b}$, $a,b\in\{0,1\}$, are stabilizer
states spanning a four-dimensional subspace of $(\C^2)^{\otimes 10}$ which contains
$\ket d\otimes\ket d$ by \eqref{eq:Chat}. Call it the \emph{tensor-square plane}
\[
\mathcal P:=\spn\{\ket{w_a}\otimes\ket{w_b}\ :\ a,b\in\{0,1\}\}.
\]
Every $\ket\psi\in\mathcal P$ has an amplitude that is constant on each of the four
\emph{cells}
\[
\mathcal C_{ab}:=\{(y,z)\in\F_2^5\times\F_2^5:\ \qq(y)=a,\ \qq(z)=b\},
\qquad |\mathcal C_{ab}|=16^2=256,
\]
and $\ket\psi$ is determined by the table $(m_{ab})_{a,b\in\F_2}$ of those values.

\begin{theorem}[Tensor-square plane rigidity]\label{thm:squareplane}
The only stabilizer rays contained in $\mathcal P$ are the four product rays
$[w_a\otimes w_b]$, $a,b\in\{0,1\}$.
\end{theorem}

\begin{proof}
Let $\ket\psi\in\mathcal P$ be a stabilizer state with cell table $(m_{ab})$. Its
support is the union of the cells with $m_{ab}\ne0$, hence has cardinality
$256\cdot\#\{(a,b):m_{ab}\ne0\}$. A stabilizer state has affine support, of cardinality
a power of $2$, so the number of nonzero cells is $1$, $2$ or $4$.

\emph{One cell.} The support is $\qq^{-1}(a)\times \qq^{-1}(b)$. The image of an affine
set under the coordinate projection $(y,z)\mapsto y$ is affine, and here that image is
$\qq^{-1}(a)$, which is not affine (Lemma~\ref{lem:qstruct}(c)). Contradiction.

\emph{Two cells.} If the two cells form a row or a column, the projection onto one of
the two blocks is again a level set of $\qq$, and the same contradiction applies. The
remaining possibility is a diagonal, i.e.\ a level set of the Boolean quadratic
$Q(y,z):=\qq(y)+\qq(z)$ on $\F_2^{10}$. Such a level set has $512=2^9$ elements, so if
it were affine it would be an affine hyperplane, and then $Q$ would be an $\F_2$-affine
function; but the polar form of $Q$ is $\bb\oplus\bb$, which is nonzero
(e.g.\ $\bb(e_1,e_2)=1$ by Lemma~\ref{lem:qstruct}(a)), while affine functions have zero
polar form. Contradiction.

\emph{Four cells (full support).} Fix $z_0\in\F_2^5$ and contract the second block with
$\bra{z_0}$. By Lemma~\ref{lem:ops}(e) the result is a stabilizer state (it is nonzero,
as $\ket\psi$ has full support); it lies in $\spn\{w_0,w_1\}$ and has full support, so
Theorem~\ref{thm:plane} (Case 1) forces its two cell values to have ratio $\pm1$. Thus
$m_{1b}/m_{0b}\in\{1,-1\}$ for $b\in\{0,1\}$, and symmetrically
$m_{a1}/m_{a0}\in\{1,-1\}$. Rescaling so that $m_{00}=1$, the table is therefore a sign
table, $m_{ab}=(-1)^{s(a,b)}$ with $s:\F_2^2\to\F_2$ and $s(0,0)=0$; write
$s(a,b)=\lambda a+\mu b+\nu ab$ with $\lambda,\mu,\nu\in\F_2$.

If $\nu=0$ the table factorizes, $m_{ab}=(-1)^{\lambda a}(-1)^{\mu b}$, and $\ket\psi$
is one of the four product rays $[w_\lambda\otimes w_\mu]$.

If $\nu=1$, the phase exponent of $\ket\psi$ is
\[
h(y,z)=2\big(\lambda\qq(y)+\mu\qq(z)+\qq(y)\qq(z)\big)\bmod4 .
\]
Reading $\qq$ as an integer-valued function with values in $\{0,1\}$, group the terms as
$h(y,z)=2\big(\lambda+\qq(z)\big)\qq(y)+2\mu\qq(z)$. Take the directions
$\mathbf a=(e_1,0)$ and $\mathbf b=(e_2,0)$, both supported in the first block, and the
base point $(0,z)$. The term $2\mu\qq(z)$ does not depend on $y$ and cancels, so
\[
(\Delta_{\mathbf a}\Delta_{\mathbf b}h)(0,z)
=2\big(\lambda+\qq(z)\big)\big[\qq(e_1+e_2)-\qq(e_1)-\qq(e_2)+\qq(0)\big]
=-2\big(\lambda+\qq(z)\big)\equiv 2\big(\lambda+\qq(z)\big)\pmod 4,
\]
since by Lemma~\ref{lem:sigma} the bracket is $1-1-1+0=-1$. As $z$ ranges over
$\F_2^5$, $\qq(z)$ takes both values, so this second difference takes the value $0$ for
some base points and $2$ for others, whichever $\lambda\in\{0,1\}$ is.
Corollary~\ref{cor:test}(b) rules this case out.
\end{proof}

\begin{corollary}\label{cor:noplane}
Any decomposition of $\ket d\otimes\ket d$ into three stabilizer states must use at
least one stabilizer state outside $\mathcal P$.
\end{corollary}

\begin{proof}
By \eqref{eq:Chat},
\[
\ket d\otimes\ket d=\frac{-i}2\Big(\ket{w_0}\otimes\ket{w_0}+i\ket{w_0}\otimes\ket{w_1}
+i\ket{w_1}\otimes\ket{w_0}-\ket{w_1}\otimes\ket{w_1}\Big),
\]
so all four coordinates of $d\otimes d$ in the basis $\{w_a\otimes w_b\}$ of
$\mathcal P$ are nonzero. By Theorem~\ref{thm:squareplane}, three stabilizer states
inside $\mathcal P$ are (up to scalars) three of the four basis vectors, and their span
is a coordinate subspace of dimension $3$, which does not contain $d\otimes d$.
\end{proof}

\subsection{One entangled slice, and the failure of its natural enlargement}

Corollary~\ref{cor:noplane} says a three-term certificate must involve a stabilizer
state entangled across the $5\,|\,5$ cut. At least one such state does interact
exactly with $d\otimes d$.

\begin{proposition}[Graph slice]\label{prop:graph}
Let $\mathcal G:=\{(y,y+\one_5):y\in\F_2^5\}\subseteq\F_2^{10}$, a linear
subspace of dimension $5$. The amplitude of $\ket d\otimes\ket d$ is the constant $-i$
on $\mathcal G$. Consequently the uniform superposition
$\ket{g}:=\sum_{y}\ket{y}\otimes\ket{y+\one_5}$ is a stabilizer state, entangled across the
$5\,|\,5$ cut, on which $d\otimes d$ restricts to $-i\ket g$.
\end{proposition}

\begin{proof}
By Lemma~\ref{lem:qstruct}(b), $\qq(y)+\qq(y+\one_5)=1$ for every $y$, so the amplitude
$i^{-\qq(y)}i^{-\qq(y+\one_5)}=i^{-1}=-i$. The set $\mathcal G$ is the graph of the
invertible affine map $y\mapsto y+\one_5$, hence a $5$-dimensional affine (indeed linear)
subspace, and a uniform superposition over an affine subspace is a stabilizer state
(Definition~\ref{def:stab} with $h=0$). It is entangled across the cut because its
reduced density matrix on the first block is proportional to the identity on a
$32$-dimensional space, while a product state would have rank one.
\end{proof}

The obvious way to try to grow this slice is to enlarge $\mathcal G$ to a subspace whose
direction space is totally isotropic for the polar form $\bb\oplus\bb$ of the exponent
of $d\otimes d$ modulo $2$. That parity condition is necessary
(Lemma~\ref{lem:diff}(i)) but not sufficient, and the natural first enlargement fails.

\begin{proposition}[The $64$-point enlargement is not stabilizer]\label{prop:noenlarge}
Let
$\mathcal L:=\{(y,y+s\one_5)\ :\ y\in\F_2^5,\ s\in\F_2\}$, a linear subspace of $\F_2^{10}$
of dimension $6$ whose direction space is totally isotropic for $\bb\oplus\bb$. The
restriction of $\ket d\otimes\ket d$ to $\mathcal L$ is not a stabilizer state.
\end{proposition}

\begin{proof}
Parametrize $\mathcal L$ by $(y,s)\in\F_2^5\times\F_2$. For $s=0$ the amplitude of
$d\otimes d$ is $i^{-2\qq(y)}=(-1)^{\qq(y)}$, and for $s=1$ it is $-i$ by
Proposition~\ref{prop:graph}. A phase exponent is therefore
\[
H(y,s)=2\qq(y)\,(1+s)+3s \bmod 4 ,
\]
since $H(y,0)=2\qq(y)$ and $H(y,1)=4\qq(y)+3\equiv3$. Taking directions $e_1,e_2$ in
the $y$ block and base point $(0,s)$,
\[
(\Delta_{e_1}\Delta_{e_2}H)(0,s)=2(1+s)\big[\qq(e_1+e_2)-\qq(e_1)-\qq(e_2)+\qq(0)\big]
=2(1+s)\cdot(-1),
\]
i.e.\ $-2\equiv2\pmod4$ for $s=0$ and $-4\equiv0\pmod4$ for $s=1$. The second difference
depends on the base point, so Corollary~\ref{cor:test}(b) applies.

For the isotropy claim: the direction space of $\mathcal L$ is spanned by
$\{(e_j,e_j)\}_{j=1}^5$ and $(0,\one_5)$. Since $\one_5$ spans the radical of $\bb$
(Lemma~\ref{lem:qstruct}(a)), the vector $(0,\one_5)$ pairs trivially with everything;
and for the remaining generators,
\[
(\bb\oplus\bb)\big((e_j,e_j),(e_k,e_k)\big)=\bb(e_j,e_k)+\bb(e_j,e_k)=0 .
\]
\end{proof}

\subsection{A finite reduction}

Finally we reduce the exclusion of $\chi(d^{\otimes2})\le3$ to a finite (if large)
amount of checking.

\begin{lemma}[No hyperplane-coset restriction is stabilizer]\label{lem:norestr}
Let $\mathcal K\subseteq\F_2^{10}$ be a coset of an affine hyperplane, i.e.\
$\mathcal K=\{\mathbf x:v\cdot \mathbf x=\epsilon\}$ for some $v\ne0$, $\epsilon\in\F_2$. Then the
restriction of $\ket d\otimes\ket d$ to $\mathcal K$ is not a stabilizer state; hence
its stabilizer rank is at least $2$.
\end{lemma}

\begin{proof}
The restriction has full support on $\mathcal K$, so if it were a stabilizer state its
affine support would be exactly $\mathcal K$, whose direction space $\mathcal K_0$ has
dimension $9$. Its phase exponent is the restriction of
$h(y,z)=-\qq(y)-\qq(z)$, and $h\equiv \qq(y)+\qq(z)\pmod2$, whose second difference in
directions $\mathbf a,\mathbf b$ is $(\bb\oplus\bb)(\mathbf a,\mathbf b)$ modulo $2$. By
Lemma~\ref{lem:diff}(i) the form $\bb\oplus\bb$ would have to vanish identically on
$\mathcal K_0$, i.e.\ $\mathcal K_0$ would be totally isotropic. But the radical of
$\bb\oplus\bb$ is $\{0,\one_5\}\times\{0,\one_5\}$, of dimension $2$
(Lemma~\ref{lem:qstruct}(a)), so by Lemma~\ref{lem:isotropy} a totally isotropic
subspace of $\F_2^{10}$ has dimension at most $2+\frac{10-2}2=6<9$. Contradiction.
\end{proof}

\begin{proposition}[Finite reduction]\label{prop:reduction}
Suppose $\ket d\otimes\ket d=\sum_{j=1}^3c_j\ket{\varphi_j}$ with
$\ket{\varphi_j}\in\Stab_{10}$ and let $\mathcal A_j$ be the affine support of
$\ket{\varphi_j}$. Then:
\begin{enumerate}[label=(\alph*)]
\item $\mathcal A_1\cup\mathcal A_2\cup\mathcal A_3=\F_2^{10}$ and at least one
$\mathcal A_j$ has dimension $\ge9$;
\item either all three $\mathcal A_j$ equal $\F_2^{10}$, or some hyperplane-coset
restriction of $\ket d\otimes\ket d$ has stabilizer rank exactly $2$;
\item if one summand is the graph slice of Proposition~\ref{prop:graph} and the other
two supports are proper, then those two supports are the two complementary cosets of
one affine hyperplane.
\end{enumerate}
Consequently $\chi(d^{\otimes2})=4$ would follow from the two finite verifications:
(i) none of the $2(2^{10}-1)=2046$ hyperplane-coset restrictions of $d\otimes d$ has
stabilizer rank $2$; and (ii) $d\otimes d$ is not a combination of at most three
full-support stabilizer states.
\end{proposition}

\begin{proof}
(a) $\ket d\otimes\ket d$ has full support, so the union of the $\mathcal A_j$ is all of
$\F_2^{10}$. If every $\mathcal A_j$ had dimension $\le8$ then
$|\bigcup_j\mathcal A_j|\le3\cdot2^8=768<1024=2^{10}$.

(b) Suppose some $\mathcal A_j$ is proper. Then it lies in an affine hyperplane
$\mathcal H$; let $\mathcal K=\F_2^{10}\setminus\mathcal H$ be the opposite coset.
Restriction to $\mathcal K$ is, up to a Clifford, the Pauli projection
$\frac12(I-(-1)^{\epsilon}Z^v)$ of Lemma~\ref{lem:ops}(d) followed by discarding the
complement; it annihilates $\ket{\varphi_j}$ and maps each remaining summand to zero or
a stabilizer state. Hence the restriction of $\ket d\otimes\ket d$ to $\mathcal K$ has
stabilizer rank at most $2$, and by Lemma~\ref{lem:norestr} it is not $1$.

(c) The graph slice has $32$ points, so the other two supports must cover the
remaining $1024-32=992$ points. Two proper affine subspaces have at most $512$ points
each; if one had dimension $\le8$ the union would have at most $512+256=768<992$
points. So both are affine hyperplane cosets. Two distinct cosets of the \emph{same}
hyperplane cover all $1024$ points, whereas two cosets of \emph{different} hyperplanes
intersect in $256$ points and cover only $512+512-256=768<992$. Hence they are the two
complementary cosets of a single hyperplane.

The final statement follows since (a)--(b) exhaust the possibilities: either all
supports are full, which is case (ii), or case (i) applies.
\end{proof}

We emphasise that neither of the two finite verifications is carried out here. Both are
large: the number of $9$-qubit stabilizer states, and a fortiori the number of pairs, is
$2^9\prod_{k=1}^{9}(2^k+1)\approx1.8\times10^{16}$, and the number of full-support
$10$-qubit stabilizer states is $4^{10}2^{45}\approx3.7\times10^{19}$. Exhausting either
would require substantially better pruning than support dimension alone.

\section{Lower bounds}\label{sec:lower}

\subsection{The bound $\chi(C\otimes C)\ge3$}

\begin{proposition}\label{prop:lower3}
$\chi(C\otimes C)\ge3$.
\end{proposition}

\begin{proof}
By Lemma~\ref{lem:catamp} the amplitude of $\ket C$ at $x=0^6$ is $\sqrt2$, so
$\braket{0^6|C}=\sqrt2\ne0$ and therefore
\begin{equation}\label{eq:contract0}
\big(\mathrm{id}\otimes\bra{0^6}\big)\big(\ket C\otimes\ket C\big)=\sqrt2\,\ket C .
\end{equation}
By Lemma~\ref{lem:ops}(e), $\chi(C)\le\chi(C\otimes C)$, and $\chi(C)=3$ by
Theorem~\ref{thm:cited}.
\end{proof}

Together with Theorem~\ref{thm:nine} this proves Theorem~\ref{thm:main}.

\subsection{A $\cat_2$ contraction and a $\cat_{12}$ projection}

The following two exact relations do not improve the numerical interval, but they
identify smaller and larger exact-rank problems whose resolution would.

\begin{proposition}[Contraction to $\cat_{10}$]\label{prop:cat10}
Let $\ket\beta:=\ket{00}+i\ket{11}$, a two-qubit stabilizer state, and contract
$\bra\beta$ against the sixth qubit of each copy of $C$. Then, up to reordering the ten
remaining qubits,
\begin{equation}\label{eq:cat10}
\big(\bra\beta\otimes\mathrm{id}\big)\big(\ket C\otimes\ket C\big)
=\ket T^{\otimes10}+\ket{T^\perp}^{\otimes10}=\sqrt2\,\ket{\cat_{10}} ,
\end{equation}
and hence $\chi(C\otimes C)\ge\chi(\cat_{10})$.
\end{proposition}

\begin{proof}
Direct computation with $\bra\beta=\bra{00}-i\bra{11}$ gives
\[
\braket{\beta|T,T}=1+(-i)\omega^2=1+(-i)(i)=2,\qquad
\braket{\beta|T^\perp,T^\perp}=1+(-i)(+1)\omega^2=2,
\]
\[
\braket{\beta|T,T^\perp}=1+(-i)(-\omega^2)=1+i\cdot i=0,\qquad
\braket{\beta|T^\perp,T}=0 .
\]
Expanding $\ket C\otimes\ket C=\frac12\big(T^{\otimes6}+ (T^\perp)^{\otimes6}\big)
\otimes\big(T^{\otimes6}+(T^\perp)^{\otimes6}\big)$ and contracting, the two mixed
terms vanish and the two like terms each acquire a factor $2$, giving
$\frac12\cdot2\,(T^{\otimes10}+(T^\perp)^{\otimes10})$, which is
$\sqrt2\ket{\cat_{10}}$ by Definition~\ref{def:T}. The rank inequality is
Lemma~\ref{lem:ops}(e).
\end{proof}

\begin{proposition}[Projection from $\cat_{12}$]\label{prop:cat12}
Let $\Pi:=\frac12\big(I+Z^{\otimes6}\otimes I^{\otimes6}\big)$, the projector onto even
parity of the first six qubits. Then
\begin{equation}\label{eq:cat12}
\ket C\otimes\ket C=\sqrt2\,\Pi\ket{\cat_{12}},
\end{equation}
and hence $\chi(C\otimes C)\le\chi(\cat_{12})$.
\end{proposition}

\begin{proof}
By Lemma~\ref{lem:catamp}, $\ket{\cat_{12}}=\sqrt2\sum_{\mathbf x\in\Ev_{12}}\omega^{\wt{\mathbf x}}\ket{\mathbf x}$.
Writing $\mathbf x=(x,x')$ with $x,x'\in\F_2^6$, the projector $\Pi$ keeps exactly those $\mathbf x$
with $\wt x$ even, and combined with $\wt x+\wt{x'}$ even this forces $\wt{x'}$ even
too. Hence
$\sqrt2\,\Pi\ket{\cat_{12}}=2\sum_{x,x'\in\Ev_6}\omega^{\wt x+\wt{x'}}\ket{x,x'}$,
which is precisely the amplitude of $\ket C\otimes\ket C$ by \eqref{eq:catamp}. The rank
inequality is Lemma~\ref{lem:ops}(d).
\end{proof}

\begin{theorem}[Sandwich]\label{thm:sandwich}
$\chi(\cat_{10})\le\chi(C\otimes C)\le\chi(\cat_{12})$.
\end{theorem}

\begin{proof}
Combine Propositions~\ref{prop:cat10} and \ref{prop:cat12}.
\end{proof}

\begin{remark}
Note that $\ket{\cat_{12}}\ne\ket C\otimes\ket C$: expanding definitions,
$\ket C\otimes\ket C=\frac1{\sqrt2}\big(\ket{\cat_{12}}+\ket{\mathrm{cross}}\big)$ where
$\ket{\mathrm{cross}}=\frac1{\sqrt2}\big(T^{\otimes6}\otimes(T^\perp)^{\otimes6}
+(T^\perp)^{\otimes6}\otimes T^{\otimes6}\big)$ is orthogonal to $\ket{\cat_{12}}$.
Consequently bounds on the tensor square do not transfer to $\cat_{12}$ and conversely;
\eqref{eq:cat12} is the only implication between them proved here. At present
\eqref{eq:cat10} does not improve the numerical lower bound: computational-basis
contraction only gives $\chi(\cat_{10})\ge\chi(\cat_6)=3$, and no proof of
$\chi(\cat_{10})\ge4$ is known to us.
\end{remark}

\section{Discussion and open problems}\label{sec:open}

\subsection{Status}

The proved numerical result is $3\le\chi(C\otimes C)\le9$
(Theorem~\ref{thm:main}). The upper bound is certified by the explicit three-term
decomposition \eqref{eq:threeterm} and its tensor square \eqref{eq:nine}; the lower
bound by $\chi(\cat_6)=3$ and the contraction \eqref{eq:contract0}. The exact relations
$\chi(\cat_{10})\le\chi(C\otimes C)\le\chi(\cat_{12})$ hold but do not currently narrow
the interval. Beyond these numbers, the structural content is:

\begin{itemize}[leftmargin=*]
\item $\chi(C\otimes C)=\chi(\widehat C^{\otimes2})$ with $\widehat C$ a full-support
five-qubit state (Theorem~\ref{thm:compress}), and
$\chi(C\otimes C)\le5+\chi(d^{\otimes2})$ with $\chi(d)=2$
(Theorem~\ref{thm:reduce}, Lemma~\ref{lem:chid}). Conjecture~\ref{conj:main} would
follow from $\chi(d^{\otimes2})\le3$.
\item An eight-term decomposition into stabilizer states that are product across the
$6\,|\,6$ cut must have dependent factor families, with local ranks $r,r'\ge3$ and
$r+r'\le9$ (Propositions~\ref{prop:indep} and \ref{prop:sylvester}).
\item The tensor-square plane $\mathcal P$ contains exactly four stabilizer rays, so no
three-term decomposition of $d^{\otimes2}$ lives inside it
(Theorem~\ref{thm:squareplane}, Corollary~\ref{cor:noplane}). Exactly one entangled
stabilizer slice is known (Proposition~\ref{prop:graph}), and its natural
maximal-isotropic enlargement is excluded (Proposition~\ref{prop:noenlarge}).
\item Excluding $\chi(d^{\otimes2})\le3$ reduces to two finite tasks
(Proposition~\ref{prop:reduction}), neither of which is completed here.
\end{itemize}

\subsection{The three remaining gaps}

\begin{enumerate}[label=\textbf{(G\arabic*)},leftmargin=*]
\item \emph{The eight-term upper bound.} What is missing is either (i) a three-term
certificate for $d^{\otimes2}$ --- necessarily using at least one stabilizer state
outside $\mathcal P$, by Corollary~\ref{cor:noplane} --- or (ii) an explicit
affine-quadratic description of eight stabilizer summands of $C\otimes C$ with an exact
amplitude verification, or (iii) the completion of the two finite tasks of
Proposition~\ref{prop:reduction}, which would instead prove $\chi(d^{\otimes2})=4$ and
close off route \eqref{eq:reduce}. Note that (iii) alone would \emph{not} refute
Conjecture~\ref{conj:main}: route \eqref{eq:reduce} is sufficient, not necessary.
\item \emph{A lower bound above three.} The contractions certify only rank three. The
$\cat_2$ contraction \eqref{eq:cat10} reduces the question to $\chi(\cat_{10})\ge4$,
for which the pair-classification technique of \cite{QPG2021} does not obviously
generalize, since it would require control of spans of \emph{three} stabilizer states.
The $\cat_{12}$ projection supplies nothing stronger. A proof of $\chi(\cat_{10})\ge4$,
or a different contraction with a known rank-four output, would improve the interval,
though it would not decide Conjecture~\ref{conj:main}.
\item \emph{The exact value.} Every value in $\{3,4,\dots,9\}$ remains compatible with
the two mechanisms used here. Determining $\chi(C\otimes C)$ requires an exact
$k$-term witness together with an exclusion of all ranks below $k$.
\end{enumerate}

\subsection{Approaches that did not work}

We record the failures so they are not repeated.

\begin{itemize}[leftmargin=*]
\item \emph{Deleting a term from \eqref{eq:nine}.} Impossible: the nine summands are
linearly independent (Theorem~\ref{thm:nine}).
\item \emph{Staying product across the cut.} Propositions~\ref{prop:indep} and
\ref{prop:sylvester} force strong dependencies but do not close the case, and they say
nothing about entangled summands.
\item \emph{Staying inside $\mathcal P$.} Excluded by Theorem~\ref{thm:squareplane}.
\item \emph{Growing the graph slice isotropically.} Excluded by
Proposition~\ref{prop:noenlarge}: total isotropy of the direction space is necessary but
not sufficient, and the third difference obstructs at the first enlargement.
\item \emph{Invoking known multiplicativity results.} \cite[Thm.~4.1]{LovitzSteffan2022}
proves $\chi(\psi_\zeta^{\otimes2})=4$ only for two-qubit $\psi_\zeta$ with $\zeta$
outside a finite exceptional set; $\ket d$ has amplitudes in $\Q(i)$ and five qubits, so
that theorem does not apply.
\item \emph{Stabilizer extent or fidelity as a rank lower bound.} The usual
Cauchy--Schwarz argument relating extent to rank degrades badly for an ill-conditioned
minimal spanning set, and we obtained no usable inequality.
\end{itemize}

Appendix~\ref{app:numerics} records the numerical searches, which produced no
certificate and are reported only as negative evidence about restricted search pools.

\subsection{Promising next steps}

For the upper bound: an exact affine-quadratic search over genuinely entangled
ten-qubit stabilizer states, seeded by Proposition~\ref{prop:graph} and pruned by
Lemma~\ref{lem:norestr} (which shows no summand can have a rank-one hyperplane-coset
restriction). For the lower bound: a classification of spans of three stabilizer states
sufficient to decide $\chi(\cat_{10})\ge4$. A cleaner intermediate target is the exact
value of $\chi(d^{\otimes2})$, which by Theorem~\ref{thm:reduce} is the crux of
route \eqref{eq:reduce} and lives on only ten qubits with full support.

\appendix

\section{Exact verification of the explicit identities}\label{app:verify}

All identities in this note involving explicit amplitudes were verified by exact
symbolic computation over $\Q(\omega)$, $\omega=e^{i\pi/4}$, namely:
Lemma~\ref{lem:catamp} (equation \eqref{eq:catamp}); Lemma~\ref{lem:sigma};
Theorem~\ref{thm:threeterm} (equation \eqref{eq:threeterm}); Lemma~\ref{lem:phiindep};
Lemma~\ref{lem:sigmaq}; Theorem~\ref{thm:compress} (equations \eqref{eq:compress} and
\eqref{eq:Chat}); Lemma~\ref{lem:qstruct}(a)--(c); the second-difference computations in
Lemma~\ref{lem:chid}, Theorem~\ref{thm:plane}, Theorem~\ref{thm:squareplane} and
Proposition~\ref{prop:noenlarge}; Corollary~\ref{cor:noplane};
Proposition~\ref{prop:graph}; and Propositions~\ref{prop:cat10} and \ref{prop:cat12}
(equations \eqref{eq:cat10} and \eqref{eq:cat12}).

In addition, Theorem~\ref{thm:plane} was independently confirmed by brute force over all
$4^5\cdot2^{10}=2^{20}$ full-support five-qubit stabilizer states: exactly two of them
lie in $\spn\{w_0,w_1\}$, namely $w_0$ and $w_1$.

\section{Numerical explorations (no certificates obtained)}\label{app:numerics}

The following searches were run as guidance. None yields a proof, and none is used
above.

\begin{enumerate}[label=(\alph*),leftmargin=*]
\item \emph{Pauli spectrum of $\cat_5$.} With $\cat_5$ normalized, the multiset of
expectation values $\braket{\cat_5|P|\cat_5}$ over the $4^5=1024$ Pauli operators is
\[
(-1/2)^{20},\ (-1/4)^{40},\ 0^{782},\ (1/4)^{120},\ (1/2)^{60},\ 1^{2}.
\]
This is the finite datum entering the rank-two exclusion of \cite{QPG2021}; on its own
it gives no control over spans of three stabilizer states and so no bound on
$\chi(\cat_{10})$.
\item \emph{Symmetry-generated pool.} The Clifford $A=e^{-i\pi/4}SX$ has $\ket T$ and
$\ket{T^\perp}$ as eigenvectors with eigenvalues $+1$ and $-1$; consequently $A^{\otimes S}$
fixes $\ket C$ for every even-cardinality subset $S$ of the six qubits. Applying these
$32$ symmetries to $\varphi_1,\varphi_2,\varphi_3$ generates exactly $48$ distinct
stabilizer rays, spanning a $32$-dimensional subspace. Randomized searches over
products of these rays found no eight-term witness. This is evidence only about that
restricted pool.
\item \emph{Simulated annealing.} Annealing at rank eight for $\widehat C^{\otimes2}$,
and at rank three for $d^{\otimes2}$ warm-started from a leave-one-out four-product
seed, produced no exact solution and did not improve on the seed residual. Failure of a
heuristic search is not evidence against Conjecture~\ref{conj:main}.
\end{enumerate}

\begin{thebibliography}{9}

\bibitem{QassimThesis}
Hammam Qassim,
\emph{Classical Simulations of Quantum Systems Using Stabilizer Decompositions},
PhD thesis, University of Waterloo, 2020 (deposited 2021).
\url{https://uwspace.uwaterloo.ca/items/147b4d1e-1a97-49a8-a3e0-f5954c570cdb}.
The decomposition giving $\chi(\cat_6(T))\le3$ is in \S4.2; Conjecture~\ref{conj:main}
is Eq.~(4.30) in \S4.4, and the conjectural exponent $3/8$ is Eq.~(4.31).

\bibitem{QPG2021}
Hammam Qassim, Hakop Pashayan, and David Gosset,
\emph{Improved upper bounds on the stabilizer rank of magic states},
Quantum \textbf{5}, 606 (2021).
\url{https://doi.org/10.22331/q-2021-12-20-606}, arXiv:2106.07740.
The appendix proves $\chi(\cat_6)\ge\chi(\cat_5)\ge3$ via a canonical form for pairs of
stabilizer states.

\bibitem{LovitzSteffan2022}
Benjamin Lovitz and Vincent Steffan,
\emph{New techniques for bounding stabilizer rank},
Quantum \textbf{6}, 692 (2022).
\url{https://doi.org/10.22331/q-2022-04-20-692}, arXiv:2110.07781.
Theorem~4.1 gives the first states with multiplicative stabilizer rank; Corollary~3.3
gives $\chi(T^{\otimes n})\ge(n+1)/(4\log_2(n+1))$.

\bibitem{PSV2022}
Shir Peleg, Amir Shpilka, and Ben Lee Volk,
\emph{Lower bounds on stabilizer rank},
Quantum \textbf{6}, 652 (2022).
\url{https://doi.org/10.22331/q-2022-02-15-652}.

\bibitem{BSS2016}
Sergey Bravyi, Graeme Smith, and John A. Smolin,
\emph{Trading classical and quantum computational resources},
Physical Review X \textbf{6}, 021043 (2016).
\url{https://doi.org/10.1103/PhysRevX.6.021043}.

\bibitem{BBCCGH2019}
Sergey Bravyi, Dan Browne, Padraic Calpin, Earl Campbell, David Gosset, and Mark Howard,
\emph{Simulation of quantum circuits by low-rank stabilizer decompositions},
Quantum \textbf{3}, 181 (2019).
\url{https://doi.org/10.22331/q-2019-09-02-181}.

\bibitem{DehaeneDeMoor2003}
Jeroen Dehaene and Bart De Moor,
\emph{Clifford group, stabilizer states, and linear and quadratic operations over GF(2)},
Physical Review A \textbf{68}, 042318 (2003).
\url{https://doi.org/10.1103/PhysRevA.68.042318}.

\bibitem{VandenNest2010}
Maarten Van den Nest,
\emph{Classical simulation of quantum computation, the Gottesman--Knill theorem, and
slightly beyond},
Quantum Information and Computation \textbf{10}, 258--271 (2010).

\end{thebibliography}

\end{document}
